\newglossaryentry{atom}{
    name={atom},
    description={A \gls{nucleus} with surrounding electron shells; makes up most matter},
    plural={atoms}
}
\newglossaryentry{nucleus}{
    name={nucleus},
    description={The core of an \gls{atom}, made up of \glspl{nucleon} (protons and neutrons)},
    plural={nuclei}
}
\newglossaryentry{electron}{
    name={electron},
    description={The first generation charged lepton.
    It is the most common lepton, and is what balances the positive charge in \glspl{atom}, forming electron shells},
    plural={electrons},
    symbol={$\electron[-]$}
}
\newglossaryentry{electron_shell}{
    name={electron shell},
    description={The different clouds of \glspl{electron} around the \gls{nucleus} of an \gls{atom}.
    Each electron shell represents a different energy level of the \glspl{electron} in them.
    An electron has the following quantum numbers: $L_\text{e}=1$, $Q=-1e$; and a mass of approximately $9.109\times10^{-31}\,$kg, which gives a rest energy of $0.511\times10^{-31}\,$MeV, or $8.187\times10^{-14}\,$J},
    plural={electrons shells}
}
\newglossaryentry{nucleon}{
    name={nucleon},
    description={The most common type of baryon, making up atomic \glspl{nucleus}.
    \Glspl{proton} and neutrons are the only nucleons.
    See section \ref{zoo:nucleons}},
    plural={nucleons},
    symbol={N}
}
\newglossaryentry{proton}{
    name={nucleon},
    description={Only stable type of baryon, making up atomic \glspl{nucleus} along with \glspl{neutron}.
    Protons are \glspl{nucleon}.
    See section \ref{zoo:nucleons}},
    plural={protons},
    symbol={$\proton$}
}
\newglossaryentry{neutron}{
    name={neutron},
    description={Makes up atomic \glspl{nucleus} along with \glspl{proton}.
    Neutrons are \glspl{nucleon}.
    See section \ref{zoo:nucleons}},
    plural={neutrons},
    symbol={$\neutron$}
}
\newglossaryentry{alpha_decay}{
    name={alpha decay},
    description={A type of radioactive decay where an unstable \gls{nucleus} releases an \gls{alpha_particle} (2 \glspl{proton} and 2 \glspl{neutron}, equivalent to a helium-2 \gls{nucleus}).
    The equation for alpha decay is:
    \begin{equation}
        \ce{\ce{_Z^AX -> _{Z-4}^{A-2}Y + ^4_2{\alphaparticle}}}
    \end{equation}
    See section \ref{alpha_decay}}
}
\newglossaryentry{alpha_particle}{
    name={alpha particle},
    description={The particle released in alpha decay. It is made up of two \glspl{proton} and two \glspl{neutron}, which is equivalent to a helium-2 \gls{nucleus}},
    plural={alpha_particles},
    symbol={$\alphaparticle$}
}
\newglossaryentry{ion}{
    name={ion},
    description={An \gls{atom} that has gained or lost an \gls{electron}, meaning that the \gls{atom} now has a non-zero charge.
    Positive ions are called cations, negative ions are called anions},
    plural={ions}
}
\newglossaryentry{beta_minus_decay}{
    name={beta minus decay},
    description={A type of radioactive decay where a \gls{nucleus} with too many \glspl{neutron} has a \gls{neutron} decay into a \gls{proton}, releasing a \gls{beta_minus_particle} (an \gls{electron}) and an electron antineutrino.
    The \gls{neutron} decaying into a \gls{proton} can be simplified to just a single down quark in the \gls{neutron} turning into an up quark in the \gls{proton}.
    The equation for beta minus decay is:
    \begin{equation}
        \ce{{\downquark} -> {\upquark} + {\betaparticle[-]} + {\antineutrino[e]}}
    \end{equation}}
    See section \ref{beta_minus}
}
\newglossaryentry{beta_plus_decay}{
    name={beta plus decay},
    description={A type of radioactive decay where a \gls{nucleus} with too many \glspl{proton} has a \gls{proton} decay into a \gls{neutron}, releasing a \gls{beta_plus_particle} (a \gls{positron}) and an electron neutrino.
    The \gls{neutron} decaying into a \gls{proton} can be simplified to just a single up quark in the \gls{proton} turning into an down quark in the \gls{neutron}.
    The equation for beta minus decay is:
    \begin{equation}
        \ce{{\downquark} -> {\upquark} + {\betaparticle[-]} + {\antineutrino[e]}}
    \end{equation}}
    See section \ref{beta_plus}
}
\newglossaryentry{beta_minus_particle}{
    name={beta minus particle},
    description={The name given to the \gls{electron} released in \gls{beta_minus_decay}},
    plural={beta minus particles},
    symbol={$\betaparticle[-]$}
}
\newglossaryentry{beta_plus_particle}{
    name={beta plus particle},
    description={The name given to the \gls{positron} released in \gls{beta_minus_decay}},
    plural={beta plus particles},
    symbol={$\betaparticle[+]$}
}
\newglossaryentry{positron}{
    name={positron},
    description={The antiparticle of an \gls{electron}.
    It has a positive charge, and can annihilate an \gls{electron}, or be produced alongside one in \gls{pair_production}},
    plural={positrons},
    symbol={$\electron[+]$}
}
\newglossaryentry{pair_production}{
    name={pair production},
    description={The process of a photon (with a minimum energy of 1.022MeV) producing an \gls{electron} and a \gls{positron}.
    See section \ref{feynman_diagrams}},
}
\newglossaryentry{neutrino}{
    name={neutrino},
    description={An uncharged lepton with an extremely small mass and high velocity.
    There are three generations of neutrino: electron neutrino ($\neutrino[e]$, generation I), muon neutrino ($\neutrino[\mu]$, generation II) and tau neutrino ($\neutrino[\tau]$, generation III).
    Each generation has more mass than the previous},
    plural={neutrinos},
    symbol={$\neutrino[]$}
}
\newglossaryentry{antineutrino}{
    name={antineutrino},
    description={An uncharged antilepton with an extremely small mass and high velocity.
    Antineutrinos are the antiparticle of \glspl{neutrino}.
    There are three generations of antineutrino, corresponding to the three generations of \gls{neutrino}: electron antineutrino ($\antineutrino[e]$, generation I), muon antineutrino ($\antineutrino[\mu]$, generation II) and tau antineutrino ($\antineutrino[\tau]$, generation III).
    Each generation has more mass than the previous},
    plural={antineutrinos},
    symbol={$\antineutrino[]$}
}
\newglossaryentry{quark}{
    name={quark},
    description={Quarks are fermions that feel the strong force.
    They are fundamental particles.
    There are six types of quark, split over three generations: up ($\upquark$, generation I), down ($\downquark$, generation I), strange ($\strangequark$, generation II), charm ($\charmquark$, generation II), top ($\topquark$, generation III) and bottom ($\bottomquark$, generation III).
    The different quarks are also split into up-type and down-type, based on their charges.
    The up-type quarks are: up, strange and top; the down-type quarks are down, charm and bottom.
    There is one up-type quark and one down-type quark per generation},
    plural={quarks},
    symbol={q}
}
% Lepton number
% Weak nuclear force
% Strong nuclear force



\newglossaryentry{UNUSED_WITHOUT_SYMBOL}{
    name={},
    description={},
    plural={}
}

\newglossaryentry{UNUSED_WITH_SYMBOL}{
    name={},
    description={},
    plural={},
    symbol={}
}
